% RESEARCH_DRAFT_SHA256: 2716a4875406d8912ad79a5f4b161dd3f40363a193e13c1088e2fc7b2354e0d2
% Layout and source-reference wrappers expanded for the web renderer.
\documentclass[11pt]{article}
\usepackage[margin=25mm]{geometry}
\usepackage{fontspec,unicode-math}
\setmainfont{Latin Modern Roman}
\setmathfont{Latin Modern Math}
\usepackage{amsmath,amssymb,mathtools,xurl,hyperref}
\hypersetup{colorlinks=true,urlcolor=blue,linkcolor=blue}
\setcounter{secnumdepth}{0}
\setlength{\parindent}{0pt}
\setlength{\parskip}{5pt}
\setlength{\emergencystretch}{3em}
\newcommand{\sref}[2]{\hyperlink{source:#1:#2}{[S#2]}}
\newcommand{\cataloguescope}[1]{\paragraph{Source record.}#1}
\newcommand{\floor}[1]{\left\lfloor#1\right\rfloor}
\begin{document}
\section*{Research-attempt provenance}
This standalone research attempt retains the complete original definition body below,
including its historical source and preparation notes. Those notes describe the
earlier catalogue preparation; the new attempt and its verification are identified
separately after that body. No original ID, catalogue status, or source statement is changed.

\section{Fast Solvers for the Electronic Schrödinger Equation}\label{statement}

\subsection{Definitions and mathematical statement}

\paragraph{Shared notation.}Standard number systems are $\mathbb N,\mathbb Z,\mathbb Q,\mathbb R,\mathbb C$. Specialized notation and index ranges follow the cited source and the definitions below.


The electronic Schrodinger equation is an eigenvalue problem $H\psi=E\psi$ for a many-electron Hamiltonian. A numerical solution usually specifies an energy or wavefunction approximation and a tolerance.
\paragraph{Record scope and qualifications.}
The catalogue title is retained as metadata. Named certificates, contracts, bounds or frameworks in that title are not asserted as results of the original source extract.
The Hamiltonian class, representation, accuracy, and baseline complexity are not specified; this remains a numerical-research direction.

Source-record limitation: Fast solution of the electronic Schrodinger equation is requested without an input class, accuracy, or complexity target.
\paragraph{Statement recorded in the supplied source.}
Develop a fast computational method for solving the electronic Schrödinger equation, with substantially better complexity than current approaches. \href{https://doi.org/10.4171/owr/2011/32}{[S2]}
\paragraph{Source status.}
The supplied catalogue labels this source record open; that metadata was checked against the supplied snapshot on 2026-10-11. The label is preserved for provenance and is not an independent verification of present open status or mathematical correctness. The catalogue literature assessment is dated 2026-08-21.
\subsection{Short English statement}
Develop faster electronic-structure solvers.
\subsection{Sources}
\begin{itemize}
\item[S1] Ulam AI and the UnsolvedMath Contributors, UnsolvedMath: A Curated Collection of Open Mathematics Problems, catalogue entry. Catalogue curation and metadata: CC BY 4.0. The source extract is retained; definitions, English summaries and qualifications are editorial additions.
\url{https://huggingface.co/datasets/ulamai/UnsolvedMath}.
\item[S2] Mathematical Methods in Quantum Chemistry (2011), Oberwolfach Reports, DOI 10.4171/owr/2011/32, pp. 1769--1843
\url{https://doi.org/10.4171/owr/2011/32}.
\item Source-specific Oberwolfach Reports licences, including CC BY-SA 4.0 where applicable, are retained; see THIRD\_PARTY\_NOTICES.md.
\end{itemize}
\paragraph{Preparation and review.}
The original source and UnsolvedMath contributors retain attribution for the source extract. Codex (GPT-6; exact serving revision not exposed) assisted with definitions, English summaries, scope notes and formatting. No new human mathematical authorship or independently verified solution is asserted. Catalogue identity checked against the complete supplied registry and current repository. Distinct upstream records, including repeated questions and broad or damaged extracts, are retained for complete coverage of the source CS category. A repeated record is not a new independent conjecture. Complete mathematical-equivalence and current-literature review has not been performed. Preparation and final review follow contribution prompts C and D, with source gaps disclosed. Structural and rendering checks do not establish a new theorem or open status.
\label{end:statement}

\section{Literature and context check}
The primary quantum-chemistry report, printed pages 1804--1805, discusses linear scaling through localized matrix structure and stresses a localization regime. The catalogue asks a much broader numerical-method question without a precision model. This attempt studies finite Hermitian discretizations and proves what a residual and a controlled matrix truncation actually certify. It does not claim that the required localization or dimension bound holds for general many-electron Hamiltonians.
\section{Research attempt: Residual certificates and a conditional truncation estimate}

\paragraph{What a small residual does certify}
Let $H$ be a finite Hermitian matrix, $v$ a unit vector, $E=v^*Hv$ and $r=(H-EI)v$. An orthonormal eigenbasis gives
\[
 \lVert r\rVert^2=\sum_j|c_j|^2(\lambda_j-E)^2.
\]
Since $\sum_j|c_j|^2=1$, some eigenvalue lies within $\lVert r\rVert$ of $E$. This proves proximity to the spectrum, not specifically to its ground energy. The matrix $\operatorname{diag}(0,1)$ with $v=(0,1)$ has residual zero and energy one although its ground energy is zero. A residual-only stopping rule therefore cannot identify the ground state without further spectral information.

Suppose in addition $\lambda_0$ is the lowest eigenvalue and every other eigenvalue is at least a certified $L>E$. The eigenvalues of $(H-\lambda_0I)(H-LI)$ are nonnegative, so
\[
 0\leq v^*(H-\lambda_0I)(H-LI)v
 =\lVert r\rVert^2-(E-\lambda_0)(L-E).
\]
The Rayleigh quotient is a convex combination of eigenvalues and hence $\lambda_0\leq E$. Combining these facts yields the explicit enclosure
\[
 E-\frac{\lVert r\rVert^2}{L-E}\leq\lambda_0\leq E.
\]
The certified lower bound on the rest of the spectrum is an additional assumption, not a number inferred from the same trial vector.

\paragraph{A conditional localization-to-truncation guarantee}
Index matrix entries by vertices of a graph with maximum degree $\Delta\geq3$, and suppose
\[
 |H_{ij}|\leq C e^{-\alpha\operatorname{dist}(i,j)},\qquad
 \alpha>\log(\Delta-1).
\]
Let $H_R$ retain entries at graph distance at most the nonnegative integer $R$ and set the others to zero symmetrically. There are at most $\Delta(\Delta-1)^{k-1}$ vertices at distance $k\geq1$ from a given vertex: enumerate paths without immediate backtracking, which include at least one path to every such vertex. With $\beta=\alpha-\log(\Delta-1)>0$, every absolute row sum of $H-H_R$ is at most
\[
 \delta_R=\frac{C\Delta}{\Delta-1}\,
       \frac{e^{-\beta(R+1)}}{1-e^{-\beta}}.
\]
For a Hermitian matrix with absolute row sums at most $\delta_R$, any eigenvector coordinate of maximal absolute value gives $|\lambda|\leq\delta_R$ from its eigenvalue equation. The operator norm is the largest absolute eigenvalue, so $\lVert H-H_R\rVert_2\leq\delta_R$. Taking minimum Rayleigh quotients in both directions proves
\[
 |\lambda_0(H)-\lambda_0(H_R)|\leq\delta_R.
\]
Each row of $H_R$ has at most $1+\Delta\sum_{k=1}^R(\Delta-1)^{k-1}$ entries. Thus a matrix-vector multiplication is linear in the matrix dimension for fixed $R,\Delta$, and the displayed error determines a sufficient $R$ for a prescribed tolerance. This counts a multiplication, not the number of eigensolver iterations or the cost of constructing the entries.

\paragraph{The unresolved complexity step}
A many-electron discretization can itself have dimension exponential in the number of physical particles or orbitals. Linear cost in that dimension is therefore not automatically a fast physical-size algorithm. Nor do Coulomb interactions and a chosen many-body basis automatically give the assumed decay with an exponent exceeding the graph-growth rate. Finally a gap certificate needed for the ground-state residual enclosure may be costly. These are concrete missing hypotheses in the localization/residual route, rather than consequences of sparse arithmetic.

\section{Verification and adversarial audit}
The residual counterexample is an exact excited eigenstate. The ground-state enclosure checks the sign using $L>E$ and requires a bound on all other eigenvalues. The truncation proof counts graph shells and states the strict decay-versus-growth condition; at equality its geometric tail bound no longer converges. Symmetric truncation preserves Hermitian structure, which is used in the norm and Rayleigh arguments.
\paragraph{Audit corrections.} Separated spectral proximity from ground-state certification, per-multiplication cost from eigensolver cost, and discretization dimension from particle count.
\paragraph{External computation.} None; symbolic arguments and boundary checks only.
\section{Scope and remaining gaps}
The broad request for a substantially faster electronic Schroedinger solver remains unresolved. The attempt provides explicit finite-dimensional certificates and a conditional sparse-multiplication bound, while identifying the missing localization, physical-size, gap and iteration controls.
\paragraph{Integration review.} Codex (GPT-6; exact serving revision not exposed) reviewed the written arguments and expanded layout and source-reference wrappers. Repository metadata and the production web renderer were checked using Python 3.13 and JavaScript with MathJax 3.2.2. This model-assisted review is not independent human proof verification; no PDF compilation or formal proof checking is asserted.

\section{Completion Estimate}
\noindent\textbf{COMPLETION: 0\%}
The calculations below the preserved source statement investigate a concrete route,
but establish no substantiated new progress toward the entire original target.
Known special cases and failed sufficient conditions are not credited as a solution.
\section{References}
The original source references and catalogue attribution are preserved in the
Sources section above. The following URLs were actually consulted for this attempt;
retention of an additional original reference does not claim fresh consultation of it.
\begin{itemize}
\item \url{https://ems.press/content/serial-article-files/46348} (consulted 2026-10-11). 
\end{itemize}
\section{Model and contributor disclosure}
The parent Codex GPT-6 model has reviewed the written mathematical arguments at their declared scopes. This is model-assisted review, not independent human verification.
OpenAI gpt-6-astra, requested reasoning configuration ultra, generated this research
attempt on 2026-10-11. The exact serving revision is not exposed. Codex GPT-6
performed orchestration and repository-format review; its exact serving revision
is also not exposed. The mathematical route, boundary checks and final adversarial
self-audit followed contribution prompts A then B in
\texttt{docs/contribution\_prompts.md}, after reading \texttt{CONTRIBUTING.md}.
Python 3.14.2 assembled the standalone documents; any mathematical computation is
specified in Verification. No theorem prover was used. No human mathematical
authorship, substantial human proof correction, or independent human proof checking
is asserted. Model-assisted review, source retrieval, and structural checks do not
establish correctness or novelty. Full private reasoning transcripts are unavailable
as reproducibility artifacts; the complete arguments, inputs and outcomes actually
used are stated here or in the companion batch audit.
\section{Final status}
Unresolved as a new research contribution. This is a concrete route investigation,
not a claim of new complete resolution or a current-status certification. A known
answer reproduced for context does not become a new contribution. The catalogue's
source-status label is preserved; this attempt status does not assert that a
source-labelled solved question is open.
\end{document}
